%% kapa: exam=1 round=30 year=2019 period=1 form=A subject=civil_law subject_ko=민법 seq=21 number=21
%% source: Q-Net 감정평가사 기출문제 (cst00309) artlSeq=5205205 file=extracted/1008/5205205/2019년 제30회 감정평가사 1차시험 문제지/2019년 제30회 감정평가사 1차시험 문제지/2019년도 제30회 감정평가사 1차 1교시 A형.hwp
%% conversion: kordoc 4.24.51 -> tools/gen1.py (auto); review=proofread-sonnet; flags=none
\begin{question}{21}
물권법정주의에 관한 설명으로 옳은 것은? (다툼이 있으면 판례에 따름)
\begin{choices}
  \item 물권은 명령이나 규칙에 의해서도 창설될 수 있다.
  \item 민법은 관습법에 의한 물권의 성립을 부정한다.
  \item 물권법정주의에 관한 규정은 강행규정이며, 이에 위반하는 법률행위는 무효이다.
  \item 대법원은 사인(私人)의 토지에 대한 관습상의 통행권을 인정하고 있다.
  \item 미등기 무허가건물의 양수인은 그 소유권이전등기를 경료하지 않더라도 그 건물에 관하여 소유권에 준하는 관습상의 물권을 가진다.
\end{choices}
\end{question}
\answer{3}
